% Part: incompleteness
% Chapter: incompleteness-provability
% Section: provability-conditions

\documentclass[../../../include/open-logic-section]{subfiles}

\begin{document}

\olfileid{inc}{inp}{prc}

\olsection{$\Th{PA}$ માટેની \usetoken{S}{derivability} શરતો}

પેનો અંકગણિત, એટલે
$\Th{PA}$, એ
$\Th{Q}$ને દરેક
!!{formula} માટેની
ગાણિતિક અનુમાનની
સ્વયંસિદ્ધિઓથી વિસ્તરાવતો
સિદ્ધાંત છે. એટલે
$\Th{Q}$માં આ આકારની
સ્વયંસિદ્ધિઓ ઉમેરીએ:
\[
(!A(0) \land \lforall[x][(!A(x) \lif !A(x'))]) \lif \lforall[x][!A(x)]
\]
દરેક !!{formula}~$!A$
માટે. નોંધો કે આ ખરેખર
એક \emph{યોજના} છે,
એટલે અનંત સ્વયંસિદ્ધિઓ
(અને $\Th{PA}$ને સાન્ત
સ્વયંસિદ્ધિ-ગણથી
વ્યાખ્યાયિત કરી શકાતું
નથી). પરંતુ ચિહ્નોની
કોઈ શ્રેણી ગાણિતિક
અનુમાનની સ્વયંસિદ્ધિનું
ઉદાહરણ છે કે નહીં તે
અસરકારક રીતે નક્કી કરી
શકાય છે; તેથી
$\Th{PA}$નો સ્વયંસિદ્ધિ-ગણ
સંગણનીય છે. $\Th{PA}$
એ~$\Th{Q}$ કરતાં ઘણો
વધુ શક્તિશાળી સિદ્ધાંત
છે. દાખલા તરીકે, સામાન્ય
રીતે ગાણિતિક અનુમાન
વાપરીને સરવાળો અને ગુણાકાર
ક્રમવિનિમયી છે એ સરળતાથી
સાબિત થાય છે. હકીકતમાં
મોટા ભાગની સાન્તવાદી
સંખ્યાસિદ્ધાંત અને
સંયોજનશાસ્ત્રની દલીલો
~$\Th{PA}$માં કરી શકાય છે.

$\Th{PA}$ સંગણનીય રીતે
સ્વયંસિદ્ધિત હોવાથી
!!{derivability} વિધેયધર્મ
$\Prf[\Th{PA}](x,y)$
સંગણનીય છે; તેથી તેનું
~$\Th{Q}$માં (અને તેથી
~$\Th{PA}$માં પણ)
નિરૂપણ થાય છે. પહેલાંની
જેમ સંબંધને નિરૂપતા
સૂત્ર માટે
$\OPrf[\Th{PA}](x,y)$
લખીશું. $\OProv[\Th{PA}](y)$
એ સૂત્ર
$\lexists[x][\OPrf[\Th{PA}](x,y)]$
હોય. તેનો સાહજિક અર્થ
``$y$ એ $\Th{PA}$ની
સ્વયંસિદ્ધિઓમાંથી
!!{derivable} છે'' થાય છે.
\sourcecorrection{OLINC-044}{સ્થિર અંગ્રેજી
મૂળમાં આ અસ્તિત્વવાચક
સૂત્રની અંદર બહારનો
સંબંધ મૂકાયો છે. અહીં
અગાઉ વ્યાખ્યાયિત
અંકગણિતીય નિરૂપક સૂત્ર
મૂક્યું છે જેથી આખી
અભિવ્યક્તિ સિદ્ધાંતની
ભાષાનું સૂત્ર રહે.}
$\Th{Q}$ની સ્વયંસિદ્ધિઓ
કરતાં થોડું વધારે
શા માટે જોઈએ?
કારણ કે વપરાતો સિદ્ધાંત
આ !!{derivability}
વિધેયધર્મ વિશેની કેટલીક
મૂળભૂત હકીકતો
!!{derive} કરી શકે એટલો
મજબૂત હોવો જોઈએ.
જરૂરી હકીકતો આ છે:
\begin{enumerate}
\item[P1.] જો
  $\Th{PA} \Proves !A$,
  તો $\Th{PA} \Proves
  \OProv[\Th{PA}](\gn{!A})$.
\item[P2.] બધાં
  !!{formula} $!A$ અને
  $!B$ માટે,
  \[
  \Th{PA} \Proves \OProv[\Th{PA}](\gn{!A \lif !B}) \lif
  (\OProv[\Th{PA}](\gn{!A}) \lif \OProv[\Th{PA}](\gn{!B})).
  \]
\item[P3.] દરેક
  !!{formula}~$!A$ માટે,
  \[
  \Th{PA} \Proves \OProv[\Th{PA}](\gn{!A})
  \lif \OProv[\Th{PA}](\gn{\OProv[\Th{PA}](\gn{!A})}).
  \]
\end{enumerate}
આ ત્રણ ગુણધર્મો સાચા
છે તે ચકાસવા માટે
!!{formula}
$\OProv[\Th{PA}](y)$નું
કાળજીથી વર્ણન કરવું અને
સંબંધિત ઔપચારિક
!!{derivation} વર્ણવવા
$\Th{PA}$ની
સ્વયંસિદ્ધિઓ વાપરવી
પડે. શરતો (૧) અને
~(૨) સરળ છે; ખરેખર
કામ શરત~(૩)માં છે.
(વિચારો કે તેમાં કેવું
કામ કરવું પડે \dots)
વિગતો આપવી કંટાળાજનક
અને અહીં રસપ્રદ નથી;
તેથી $\Th{PA}$ ઉપરના
ત્રણ ગુણધર્મો ધરાવે છે
એ વાત અહીં માન્ય રાખવા
કહીએ છીએ. $\OProv[\Th{PA}](y)$ની
યોગ્ય પસંદગી આ પણ
સંતોષે છે:
\begin{enumerate}
\item[P4.] જો
  $\Th{PA} \Proves \OProv[\Th{PA}](\gn{!A})$,
  તો $\Th{PA} \Proves !A$.
\end{enumerate}
પરંતુ આ હકીકતની
આપણને જરૂર પડશે નહીં.

\begin{digress}
એક વાત વધુ: ગોડેલ પણ
અત્યારે આપણે છીએ એટલા
જ વિગતોમાં આળસુ હતા.
૧૯૩૧ના લેખના અંતે તેમણે
બીજા અપૂર્ણતા પ્રમેયની
સાબિતીની રૂપરેખા આપી
અને પછીના લેખમાં વિગતો
આપવાનું વચન આપ્યું.
પરંતુ તેમણે એ વિગતો
ક્યારેય આપી નહીં;
દલીલ સમજતા સૌ માનતા
હતા કે તેને પૂરું કરી
શકાય છે (એટલે તેમને
વિગતો ભરવાની જરૂર
પડી નહીં).
\end{digress}

\end{document}
